\begin{frame}
    \frametitle{Alternative Models}
    \centering
    \begin{tabular}{@{}lrrr@{}}
        \toprule
         & Log-likelihood & AIC & Units lost \\
        \midrule
        One jump for every developer & \SingleLogLikelihood & \SingleAIC & \SingleUnitsLost \\
        Non-optimizers & \NonOptimizerLogLikelihood & \NonOptimizerAIC & \NonOptimizerUnitsLost \\
        Main model: jump varies by developer & \MainLogLikelihood & \MainAIC & \MainUnitsLost \\
        Splitting cost by lot area & \LotSplittingLogLikelihood & \LotSplittingAIC & \LotSplittingUnitsLost \\
        Splitting cost by project size & \SizeSplittingLogLikelihood & \SizeSplittingAIC & \SizeSplittingUnitsLost \\
        \bottomrule
    \end{tabular}
\end{frame}

% Estimates with a kink are from the version merged into main as 6e52ba6,
% whose model allowed a cost rising above 100 units; the current code models
% the jump only.
\begin{frame}[label=kink-appendix]
    \frametitle{Allowing a Cost That Rises Above 100 Units}
    The burden becomes $b_i\big[\kappa+\tau\big((n/99)^2-(100/99)^2\big)\big]$ for a building of $n\ge100$ units.
    \par\vspace{0.8em}
    \centering
    \begin{tabular}{@{}lrrr@{}}
        \toprule
         & Kink $\tau$ & \multicolumn{2}{c}{Units lost} \\
         & with a kink & with a kink & jump only \\
        \midrule
        Main model: jump varies by developer & 0 (0--0.10) & 761 & \MainUnitsLost \\
        One jump for every developer & 0.05 & 1,130 & \SingleUnitsLost \\
        Least squares, one jump & 0.05 & 1,018 & \LeastSquaresUnitsLost \\
        \bottomrule
    \end{tabular}
    \par\vspace{0.8em}
    \raggedright
    {\footnotesize The kink is zero in the main model and matters only when every developer faces the same jump.}
    \vfill\hfill\hyperlink{main-estimates}{\beamergotobutton{Back}}
\end{frame}
